\section{Conclusion}
\label{sec:conclusion}

This work introduced RA-RMPPI, a force-conditioned TD-MPC2 controller that
adapts both its planning objective and its computation to causal contact
regimes. Matched component studies separated objective adaptation from compute
adaptation, while behaviour-cloning, router, and risk-representation ablations
identified the smallest supported formulation. In the factor-separated
simulation study, RA-RMPPI raised compound pass from 85/135 to 100/135, reduced
MRE and NRMSE, and used fewer planning samples and less host planning time
than the fixed planner built from the same checkpoints.

The improvement is specific to force regulation and computation. Task
completion remained 132/135 versus 133/135, both planners recorded two native
samples above 15~N, and the adaptive planner had a higher mean peak force.
The cylindrical 12~N condition therefore remains a concrete boundary rather
than an exception hidden by aggregate results.  A 420-evaluation
deployment-gap panel further identified signed force bias, 50-ms delay, and
registered-normal error as distinct sensitivities. In a zero-shot MuJoCo
stress test, all 30 flat and inclined evaluations retained the compound
outcome, while all 15 weak-curvature evaluations failed; the port's remaining
tangential calibration mismatch precludes a matched-solver claim. Within the
evaluated simulation setting, the study shows that continuous regime adaptation can
make direct learned wiping more accurate and less computationally demanding,
while reliable force-tail control under curvature requires a surface-following
posture or learned-constraint extension beyond the present translation-only
action interface, followed by physical validation.
